% ============================================================
% 04_siguientes_pasos.tex — Siguientes pasos
% ============================================================

\section{Siguientes Pasos}

Para la semana 6 tenemos planeado:

\begin{enumerate}
    \item \textbf{Implementar la lógica del simulador}: completar
          \texttt{c\_lock\_manager.c} y \texttt{core.c} con la lógica real de
          sincronización y clock-gating, y validar que el comportamiento sea
          correcto antes de medir nada.

    \item \textbf{Agregar escenario spin-lock}: una vez que C-Lock funcione,
          implementar un escenario equivalente con spin-lock clásico para poder
          comparar el consumo de energía entre ambos mecanismos.

    \item \textbf{Arrancar con el análisis de Selfie}: necesitamos revisar el
          código de Selfie para entender qué partes habría que modificar para
          soportar C-Lock — o al menos una versión simplificada. Esto lo
          dejamos para la semana 6 porque primero queríamos tener claro el
          paper.

    \item \textbf{Redactar el informe final}: con la base que tenemos en
          \texttt{latex/sections/}, ir completando las secciones de arquitectura
          y simulación a medida que avancemos con el código.
\end{enumerate}
